\section{The Propers: Detailed Commentary}

\subsection{A promise answered by hearing (Introit)}

The Introit antiphon puts God's promise first. He is the salvation of his
people; he hears their cry in every tribulation; his lordship endures.
The promise reaches beyond relief from one emergency. A distress ends,
but the people remain his people. Their relationship with the Lord is
larger than the event that first makes them call upon him. The attached
psalm verse then gives the people their first action: attend and incline
the ear. They enter worship as those addressed, before they become those
who offer a reply.

Psalm 77(78), of which only verse 1 is appointed, teaches through Israel's
remembered history. Its fuller course includes deliverance, nourishment,
unbelief, mercy, and the shepherding of David. Food can be received without
trust; lips can flatter while the heart remains unsteadfast. The remembered
gifts expose the inadequacy of possessing benefits without answering the
giver. The psalm therefore makes hearing more than the acquisition of
information. A people must learn to live from what God has done for them.
Its treatment of nourishment supplies a searching background to a Gospel
whose guests are summoned to food already prepared.

Augustine begins his exposition by addressing the Christian congregation
through this history. God's gifts were common, but their fruit was not
identical in every recipient. He also acknowledges grace among the faithful
of the ancient people. Their story cannot become a flattering contrast
between a wholly faithless Israel and Christians guaranteed to respond
well. It addresses Christians precisely because they too can receive much
and answer poorly.\footnote{Augustine, \work{Enarrationes in Psalmos} 77,
\S\S1--3; NPNF I.8, English heading Psalm 78. Context: Ps 77(78):19--38,
70--72.}

Augustine hears God speaking through the prophet in the opening summons.
Bellarmine identifies the immediate speaker as David, arguing from the
subsequent reference to what the fathers have told. Their explanations of
the voice differ. Both nevertheless direct the hearer toward the authority
of divine instruction rather than the prestige of the human speaker.
Neither explanation turns the psalm verse into a biblical address for the
antiphon: the Missal gives no such address for \textit{Salus populi}.
\footnote{Bellarmine, \work{Commentary on the Psalms}, Ps 77:1--8;
Augustine, \work{Enarrationes} 77.1--3.}

\subsection{Freedom for service (Collect)}

The Collect asks that everything adverse be removed, but its purpose clause
governs the request. Freedom is sought so that the worshippers may perform
what belongs to God. Adversity is not simply whatever inconveniences a
private plan, and freedom is not the ability to pursue every preference
without resistance. The prayer joins body and mind in an ordered service.
Its repeated language of being unimpeded makes obstruction the immediate
problem and readiness for God the desired good.

Schuster develops this ordering by describing the body's service under the
soul and the soul's elevation toward God through grace. Bodily health and
strength have a purpose; they are not despised, nor made the final end.
The whole person is directed toward God. His account thus interprets the
Collect's freedom as harmony, not emancipation from every bond.
\footnote{Blessed Ildefonso Schuster, \work{The Sacramentary}, vol.~III
(English 1927), ``Nineteenth Sunday after Pentecost,'' pp.~171--172.}

The Offertory later places the worshipper in the midst of tribulation.
Together the prayers allow petition for deliverance and trust during a trial
that continues. The Collect does not promise that faithful people will never
become ill, meet opposition, or remain dependent on help. It asks that their
whole life be available to God. That petition gives the Epistle's demanding
commands their proper setting: the Church asks for the freedom to obey
before she congratulates herself on any achievement.

\subsection{The new person in a common body (Epistle)}

Ephesians 4 first names the one body and its common life, then contrasts
the old manner of living with renewal in Christ. The appointed passage
begins at verse 23, with renewal in the spirit of the mind, and ends at
verse 28, with giving to the needy. The surrounding verses explain that
movement, but its appointed ending deserves particular attention. Paul
does not stop at a private conscience cleared of theft. Work reaches
beyond the worker to the person who needs help.

Chrysostom observes that renewal changes the same person, rather than
replacing human nature with another nature. The old and new person name
opposed ways of living. His clothing image therefore expresses something
more substantial than appearance and more hopeful than an unchangeable
identity as a wrongdoer. A person who has lived one way can be renewed.
At the same time, Paul's language of creation keeps this renewal dependent
on God's action. The new garment received in baptism calls for a life
answering the gift.\footnote{John Chrysostom, \work{Homilies on Ephesians}
XIII, on Eph 4:22--24; NPNF I.13.}

The transition to particular actions is essential. Chrysostom says that
particulars make the general description easier to learn. Truth is required
because members belong to one another. He illustrates the claim with the
body's mutual services: an eye that concealed danger from a foot would
endanger the body to which it too belongs. Lying is not only a private
blemish in the liar. It corrupts the means by which a neighbour can act
safely. Mutual dependence supplies a reason for truthful speech deeper
than a calculation of whether deceit will be detected.

Anger then receives an immediate boundary. The sun's setting refuses its
indefinite cultivation. Chrysostom describes the way solitude at night
can feed an injury rehearsed during the day, and how hostility then makes
false accusations seem persuasive. Resentment damages judgment as well
as affection. His insistence on prompt reconciliation answers this growth:
delay does not leave anger unchanged. The appointed injunction to give
no place to the devil concerns the breach opened within a body whose
members turn against one another.\footnote{Chrysostom, \work{Ephesians}
XIV, on 4:25--27. His account of the eye and foot and of anger's overnight
growth belongs to the homily's direct exposition of these verses.}

Finally, the hands must do something good. The cessation of stealing is
necessary, but the new person has a positive vocation to work and share.
Chrysostom distinguishes honest labour from activity that remains harmful:
the thief also exerts effort. The moral question concerns both what work
does and whom its fruits can serve. The needy person is not merely an
occasion for the worker to feel virtuous; his need determines the purpose
Paul explicitly gives to the labour.\footnote{Chrysostom, \work{Ephesians}
XIV, on 4:28. The appointed Epistle stops before the homily's ensuing
exposition of speech in 4:29.}

\subsection{Prayer before God's sight (Gradual)}

The Gradual divides the two clauses of a single verse. Prayer is compared
to incense, and uplifted hands to evening sacrifice. Both ask for an
offering to be received before God. Their force is relational: neither
the beauty of incense nor the visible gesture can make acceptance a
possession that the worshipper controls. The psalm petitions for prayer
itself to be rightly directed.

Augustine reads the evening sacrifice in Christ's voluntary Passion.
Christ offers himself; the Church prays as his body, in union with her
Head. The image thus gathers the worshipper's small act into Christ's
offering. Prayer does not become a second, independent sacrifice capable
of adding what his Passion lacks. It participates in the life of the one
whose self-gift brings his members to God.\footnote{Augustine,
\work{Enarrationes} 140.2--4; NPNF I.8, English heading Psalm 141.}

Bellarmine follows the petition into the conditions of prayer. Pure
intention directs it toward God; desire for admiration pulls it sideways.
Christ's priesthood and divine help make its ascent possible. The
surrounding psalm's concern for the mouth then shows why worship and
ordinary speech cannot be separated. The lips that pray need guarding
against excuses and harmful words.\footnote{Bellarmine, \work{Psalms}
140:1--5. The surrounding verses are context, not additional appointed
verses.}

Evening belongs to the scriptural image and to its Christological
reception. It does not prescribe the hour of this Sunday's Mass. Nor does
the mention of raised hands itself direct a particular ceremonial gesture
at the Gradual. The chant gives words to a petition for accepted prayer;
its literal and spiritual meanings remain intelligible whenever the Mass
is celebrated.

\subsection{Praise becomes proclamation (Alleluia)}

Psalm 104(105):1 joins three actions: giving glory, invoking the Lord's
name, and making his deeds known among the nations. The complete psalm
remembers covenant, Joseph, the Exodus, nourishment, and the gift of land.
It ends with the people's observance of God's statutes. The announced
deeds therefore have content. Praise tells what God has done and draws
hearers into the grateful obedience those deeds deserve.

Augustine understands \textit{confitemini} here as praise, rather than
confession of sins. He reads the proclamation among the nations as
prophetic of the evangelists. Bellarmine explains the same outward
movement through love: the lover wishes the beloved to be known by
others. He also makes invocation necessary even for praise, because the
worshipper needs God's help to praise well.\footnote{Augustine,
\work{Enarrationes} 104.1--3, English heading Psalm 105; Bellarmine,
\work{Psalms} 104:1--2.}

This outward speech follows an Epistle that has forbidden deceit.
The same community that announces God's deeds must speak truthfully
within its own life. The Alleluia's joy is generous rather than
possessive: God's goodness is proclaimed so that others may know him.
The royal servants in the Gospel will likewise go out with an invitation
whose purpose is a full feast.

\subsection{A feast both open and searching (Gospel)}

Matthew places this parable within the Temple controversy. The chief
priests and Pharisees have recognized themselves in the preceding
parables; Jesus continues speaking to them. The first refusal treats the
king's invitation as an interruption. Some go to field or business, while
others assault and kill the servants. The narrative distinguishes
indifference and murderous rejection without making either an adequate
answer to the invitation. The feast's goodness does not compel its
acceptance.

The king's reaction is severe. He judges the murderers and burns their
city, then sends the servants to gather guests from the roads. The
unworthiness of those first invited does not empty the feast. Yet the
new guests are expressly bad and good. Their gathering is followed by
the king's inspection and the expulsion of the guest without the wedding
garment. The parable therefore turns its scrutiny from refusal outside
to presence inside. Any hearer who regards himself as safely included
has still to hear its final scene.

Augustine and Gregory both interpret that mixed gathering as the present
Church. The Church's visible assembly is not identical with the company
in which evil has finally disappeared. Gregory allows more than one
account of the relation between Matthew's banquet and Luke's; his own
preferred distinction emphasizes the present mixed feast and the final
banquet from which no admitted guest departs. The moral demand does
not depend on making every narrative detail in the two accounts
identical.\footnote{Augustine, Sermon 90.1--4, NPNF I.6, Sermon XL;
Gregory the Great, \work{Homilies on the Gospels} II.38.1--3, 7.}

Their garment is charity. Augustine reaches that conclusion by asking
which good distinguishes fruitful recipients from people who possess
great gifts without living well. Gregory connects it to the Bridegroom's
own coming in love. The king's inspection consequently concerns the
heart's habit, not a modern social code about expensive clothing. Matthew
does not explain the guest's lack through a universal custom of hosts
supplying garments, and that proposed custom supplies no premise here.

The ending makes judgment real. Augustine understands the single guest
as representing a class; he does not calculate a salvation rate from one
person expelled out of a crowded hall. Gregory addresses the uncertainty
of one's own perseverance. The hearer is summoned to conversion, not
given a census of the final kingdom. Neither the violence within the
parable nor the king's judgment authorizes the congregation to retaliate
against its opponents. The scrutiny reaches the congregation itself.
\footnote{Augustine, Sermon 90.4--10; Gregory, \work{Gospel Homilies}
38.9--14.}

\subsection{Life amid tribulation (Offertory)}

The Offertory's confidence begins inside trouble. The speaker walks in
tribulation and faces hostile anger; God's hand extends to save.
Psalm 137(138), in full, gives thanks for answered prayer, celebrates
God's regard for the lowly, and trusts him to complete his work.
Thanksgiving and need coexist. Gratitude for what God has already done
does not remove the necessity of relying on him again.

Augustine presses the meaning of secure life beyond temporal immunity.
The enemy's power is real but limited: he cannot take the life that God
keeps. At the psalm's conclusion Augustine also rejects reliance on
one's own works as though the creature could complete himself apart
from the Creator.\footnote{Augustine, \work{Enarrationes} 137.10--13,
English heading Psalm 138. Verse 8 belongs to the psalm's context, not
to this appointed antiphon.}

This confidence can sustain charity toward enemies without calling
enmity harmless. The king's guests need not master every threat before
they can love. The saving right hand belongs to God, whereas the hands
named in the Epistle belong to workers who are commanded to share.
Trust and action meet without becoming identical agencies. Human
generosity is a response made possible within a life upheld by God.

\subsection{The gift and its saving fruit (Secret)}

The Secret names the gifts offered before divine majesty and asks that
they be saving \emph{for us}. Its brevity concentrates attention on the
participants. An offering's sacred character does not make their need
for prayer disappear. They ask to benefit from the gift rather than
regarding their act of offering as proof that the desired benefit is
already complete.

Schuster explains the distinction sacramentally. The Eucharist contains
the source of grace and is in itself a pledge of salvation, while its
fruit in the recipient concerns that person's dispositions. The prayer
asks for dispositions through which the oblation's saving power may
work without obstruction.\footnote{Schuster, \work{Sacramentary} III,
``Nineteenth Sunday after Pentecost,'' p.~173.}

That account gives precise force to the Gospel's question about the
guest within. It neither makes the sacrament dependent for its reality
on a recipient's goodness nor treats disposition as irrelevant. The
assembly stands before majesty in need of mercy. Charity's garment
cannot be supplied by pride in the fact that one has reached the altar.

\subsection{A command answered by desire (Communion)}

Psalm 118(119):4--5 moves abruptly from God's command to the speaker's
wish. God has commanded diligent observance; the speaker longs for
his own ways to be directed accordingly. Augustine explains
\textit{nimis} as great diligence, not excessive obedience, and
\textit{utinam} as petition. Knowing a command does not by itself
establish the strength to fulfil it. The will asks God for what he
commands.\footnote{Augustine, \work{Enarrationes} 118, opening Aleph
exposition, \S\S4--5; NPNF I.8, Psalm 119.}

Bellarmine places the desire within the relation of grace and obedient
growth. Initial justification is not earned by law; the justified
person nevertheless really grows through obedience enabled by grace.
The command is therefore neither a ladder climbed without help nor a
demand that can be discarded once grace is named.\footnote{Bellarmine,
\work{Psalms} 118:4--8.}

Schuster emphasizes the joy of the wish, its delight in God's service.
Augustine emphasizes dependence. The same wish can contain both:
the speaker loves the path and needs to be guided upon it. In the
Communion's position this is a particularly searching form of speech.
The worshipper receives and still asks to become obedient. Reception
does not turn longing into self-certification.

\subsection{Medicine for enduring attachment (Postcommunion)}

The Postcommunion asks God's healing operation to do two things:
mercifully free the worshippers from their perversities and make them
always adhere to his commands. Deliverance has a positive purpose.
Freedom from what distorts a life opens the way to a lasting attachment
to God. The prayer's \textit{semper} reaches beyond the end of this
celebration into the life that follows it.

Schuster interprets this medicine as Eucharistic grace healing desire.
The problem is not only a deficit of information about what one ought
to do. The person can know the good and remain drawn elsewhere.
Healing must reach that inclination, so that obedience can become a
living adherence rather than reluctant outward conformity.
\footnote{Schuster, \work{Sacramentary} III, pp.~173--174.}

The Collect's unimpeded service returns here as liberation from
perversity. What was requested at the opening is sought again after
Communion in a medicinal form. The new person of the Epistle and the
garment of the Gospel belong to this continuing conversion. The
congregation leaves with a petition for perseverance, not a declaration
that its members no longer need the physician.
